\begin{frame}[fragile]{A read hoisted $d$ steps buys latency hiding and owes three things back}

Issued naively, every read is followed by a wait for its own completion: the machine stalls the full
local-memory latency at each step. Hoisting the read $d$ steps creates three obligations --- and each
is a \key{closed form}, never a simulation.

\begin{center}
\begin{tikzpicture}[font=\scriptsize]
  \node[gbox] (cap) {\textbf{1. Capacity}\\$W \ge d+1$ live generations};
  \node[gbox,right=5mm of cap] (pro) {\textbf{2. Prologue}\\prime steps $0\dots d{-}1$};
  \node[gbox,right=5mm of pro] (drn) {\textbf{3. Drain}\\suppress trailing reads};
\end{tikzpicture}
\end{center}

\begin{lstlisting}[style=pseudo]
dr   = off(read, substep)                 # the REQUEST, per op-class
dr_g = 0  if fan(op) > 1                                        [veto 1]
       0  if a free sibling is outer to the group's reduction   [veto 2]
       max(0, min(dr, W_g - 1))                    # the capacity clamp
\end{lstlisting}

\begin{center}\scriptsize
\begin{tabular}{lll}
\toprule
 & \good{Shape A} ($dr_g \ge 1$) & \bad{Shape B} ($dr_g = 0$) \\
\midrule
read placement & hoisted ahead of its consumer & in place, with a wait \\
prologue       & non-empty & \textbf{empty} \\
steady advance & non-zero  & \textbf{0} \\
\bottomrule
\end{tabular}
\end{center}

\tsub{Mixing them is what the \texttt{dr\_g} refactor exists to prevent --- and was the decoder's real state: a Shape-A prologue with a Shape-B steady primes the pipe once and never re-primes it.}
\end{frame}
